\section{Introduction}

Language-model programming changed in nature within a few years. What began as code suggestions drawn from isolated prompts grew into systems that split the development cycle across specialized agents, keep persistent artifacts, run tools, and validate their own work. ChatDev and MetaGPT showed that organizing agents around engineering roles, rather than merely improving the prompt, changes what a system can deliver \cite{ChenQian2023,SiruiHong2023}. The more recent literature already speaks of AI-native software engineering and treats multi-agent systems as a research line in their own right \cite{hou2024llmse,AhmedEHassan2024,JundaHe2025}. The unit of work, once the prompt, has become the process.

This shift brings a new object of study: the operational framework that organizes development on top of an agent. Characterizing it means examining the function of each mechanism, how the specification drives generation, how repository context is retrieved, how roles are distributed, how execution acts on the environment, and how validation anchors the result, instead of surface metrics such as command counts or popularity. Existing reviews approach the problem from a different angle. Some classify agent architectures abstractly, by cognitive component; others map language-model techniques across isolated lifecycle tasks. What is absent is a software engineering lens trained on the frameworks that organize the development cycle at the level of process and repository, and a confrontation of their structure with the available evidence rather than a study-by-study description.

To fill that gap, we conducted a systematic literature review under the Kitchenham protocol, restricted to academic papers and archived preprints with verifiable records \cite{kitchenham2009slr}. From a final corpus of 37 studies published between 2022 and 2026, we propose a six-dimension taxonomy for classifying these frameworks and validate it empirically against the corpus through a hybrid thematic synthesis. The article makes three contributions: a six-dimension taxonomy (specification, context, roles, execution, validation, and portability) tested against the evidence rather than assumed a priori; a thematic synthesis of six inductive themes that covers 100\% of the corpus and exposes three axes of tension and two emergent axes outside the taxonomy; and a process-oriented research agenda that follows directly from the observed gaps.

The review is guided by five research questions. RQ1: which architectural dimensions distinguish these frameworks? RQ2: how do they organize context, roles, artifacts, execution, and validation? RQ3: what kinds of evidence support their claims? RQ4: which risks and limitations are reported or inferable? RQ5: which research gaps remain for comparing, evaluating, and governing these frameworks in real projects? The rest of the article is organized as follows. The Review Method describes the protocol and the selection flow; Related Work positions the review against the existing literature; the Proposed Taxonomy presents the six dimensions; Results and Synthesis consolidate the six themes and validate the taxonomy; the Discussion lays out the transversal findings and the axes of tension; and the Research Agenda and Conclusion outline the next steps.
